% =====================================================================
%  PRIMER PART I -- THE INVARIANCE CLUSTER
%  Four background sections for a strong undergraduate:
%    (1) orthogonal projectors,
%    (2) the discrete Fourier transform and circulant/shift structure,
%    (3) group actions, invariance/equivariance, group averaging,
%    (4) shift-invariant descriptors: power spectrum and bispectrum.
%
%  This fragment supplies NO preamble: the assembler provides
%  amsmath/amsthm and the theorem-like environments
%  (definition, theorem, lemma, corollary, proposition, example,
%   problem, remark) and the macros \R \Z \one \sign \conv \Orb.
%
%  Notation conventions are FIXED here, once, to the engineering/forward
%  DFT convention of the project digest (see the "Notation" box below):
%      omega = e^{-2 pi i / N},  F_{jk} = omega^{jk},  DFT(x) = F x.
%  All worked numbers were checked by hand and by computer.
% =====================================================================

\section{Orthogonal projectors: the one linear-algebra tool everything rests on}
\label{sec:projectors}

Almost every later idea in this primer turns out to be the same gesture: take an
object, throw away the part of it that some symmetry can move around, and keep only
the part the symmetry leaves alone. The mathematical name for ``keep only the part a
subspace can see'' is \emph{orthogonal projection}. We start here so that the rest of
the document can lean on one clean theorem instead of re-deriving it three times.

We work in a real or complex inner-product space $V$ (think of $V = \R^n$ or $V =
\mathbb{C}^n$ with the usual dot product). The \emph{inner product} of two vectors is
written $\langle u, v\rangle$; in $\R^n$ it is $\langle u,v\rangle = \sum_k u_k v_k$,
and in $\mathbb{C}^n$ it is $\langle u,v\rangle = \sum_k u_k \overline{v_k}$, where
$\overline{c}$ is the \emph{complex conjugate} of $c$ (if $c = a + bi$ then
$\overline{c} = a - bi$). The \emph{norm} is $\|v\| = \sqrt{\langle v,v\rangle}$, and
two vectors are \emph{orthogonal}, written $u \perp v$, when $\langle u,v\rangle = 0$.
Throughout, an \emph{operator} means a linear map $P : V \to V$, which we may identify
with a square matrix once a basis is fixed.

\subsection{Two adjectives: idempotent and self-adjoint}

\begin{definition}[Idempotent]
\label{def:idempotent}
An operator $P : V \to V$ is \emph{idempotent} if $P^2 = P$, i.e.\ $P(Pv) = Pv$ for
every $v \in V$. In words: applying $P$ twice does exactly what applying it once does.
\end{definition}

The intuition is that $P$ \emph{lands} every vector somewhere, and once a vector has
already landed, $P$ leaves it alone. The set of places vectors can land is the
\emph{range} (or image) $\operatorname{range}(P) = \{Pv : v \in V\}$; the set of
vectors $P$ sends to $0$ is the \emph{kernel} (or null space) $\ker(P) = \{v : Pv =
0\}$. Idempotence says precisely that $P$ acts as the identity on its own range: if
$w = Pv$ is in the range, then $Pw = P(Pv) = Pv = w$.

\begin{example}[A $2\times 2$ idempotent that is \emph{not} a clean projection]
\label{ex:oblique}
Take, for a fixed number $\alpha \neq 0$,
\[
  P \;=\; \begin{bmatrix} 0 & 0 \\ \alpha & 1 \end{bmatrix}.
\]
A direct multiplication gives
\[
  P^2 \;=\; \begin{bmatrix} 0 & 0 \\ \alpha & 1 \end{bmatrix}
            \begin{bmatrix} 0 & 0 \\ \alpha & 1 \end{bmatrix}
        \;=\; \begin{bmatrix} 0\cdot 0 + 0\cdot\alpha & 0\cdot 0 + 0\cdot 1 \\
                              \alpha\cdot 0 + 1\cdot\alpha & \alpha\cdot 0 + 1\cdot 1
              \end{bmatrix}
        \;=\; \begin{bmatrix} 0 & 0 \\ \alpha & 1 \end{bmatrix} \;=\; P,
\]
so $P$ is idempotent: it \emph{is} a projection. But it is a \emph{skewed} one. Its
range is the line spanned by $(0,1)$ and its kernel is the line spanned by
$(1,-\alpha)$ (check: $P\,(1,-\alpha)^\top = (0, \alpha - \alpha)^\top = 0$). These two
lines are \emph{not} perpendicular when $\alpha \neq 0$. Such a $P$ is called an
\emph{oblique} projection. This example, due to Wikipedia's projection article
\citep{wiki_projection}, is the cautionary tale of the section: idempotence alone does
not make a projection ``drop straight down.'' One more property is needed.
\end{example}

\begin{definition}[Self-adjoint]
\label{def:selfadjoint}
An operator $P : V \to V$ is \emph{self-adjoint} if
\[
  \langle Pu, v \rangle \;=\; \langle u, Pv \rangle \qquad \text{for all } u,v \in V.
\]
In matrix terms this is $P = P^\top$ for real $V$ (the matrix equals its transpose)
and $P = P^\ast$ for complex $V$, where $P^\ast = \overline{P}^{\top}$ is the
\emph{conjugate transpose} (transpose, then conjugate every entry). Self-adjoint real
matrices are also called \emph{symmetric}.
\end{definition}

\begin{example}[The clean version of the same idea]
\label{ex:lineproj}
Let $a = (1,1)^\top$ and form
\[
  P \;=\; \frac{a\,a^\top}{a^\top a}
        \;=\; \frac{1}{2}\begin{bmatrix} 1 & 1 \\ 1 & 1 \end{bmatrix}.
\]
Then $P = P^\top$ (it is visibly symmetric), and $P^2 = \tfrac14
\left[\begin{smallmatrix} 2 & 2 \\ 2 & 2 \end{smallmatrix}\right] = P$. So $P$ is
\emph{both} idempotent and self-adjoint. Geometrically it projects every vector
straight down onto the line through $a$: for $x = (3,1)^\top$ we get $Px = (2,2)^\top$,
and the leftover piece $x - Px = (1,-1)^\top$ is orthogonal to $a$ (their dot product
is $1 - 1 = 0$). That is the behaviour we want, and the next theorem says it always
happens \citep{projectors_tamu,wiki_projection}.
\end{example}

\subsection{The central theorem}

\begin{theorem}[Idempotent {\boldmath$+$} self-adjoint {\boldmath$\Rightarrow$}
orthogonal projection]
\label{thm:orthoproj}
Let $P : V \to V$ be idempotent ($P^2 = P$) and self-adjoint ($P = P^\ast$). Then $P$
is the \emph{orthogonal} projection onto its range $W = \operatorname{range}(P)$:
for every $v$, the vector $Pv$ lies in $W$, the leftover $v - Pv$ is orthogonal to
all of $W$, and $Pv$ is the unique closest point of $W$ to $v$.
\end{theorem}

\begin{proof}
The whole proof rests on showing one orthogonality statement; everything else follows
from it. We follow Garrett's argument \citep{projectors_garrett}.

\emph{Step 1: the complementary operator.} Write $I$ for the identity and set $Q = I -
P$. Then $Q$ measures the leftover: $Qv = v - Pv$. Note $Q$ is also self-adjoint,
because $Q^\ast = (I - P)^\ast = I^\ast - P^\ast = I - P = Q$, using $P^\ast = P$.

\emph{Step 2: the key product vanishes.} Compute, for the operator $QP = (I-P)P$,
\begin{align*}
  QP \;&=\; (I - P)P \\
     \;&=\; IP - P^2 \\
     \;&=\; P - P \qquad\text{(here we used idempotence $P^2 = P$)} \\
     \;&=\; 0 .
\end{align*}
So $QP = 0$ as an operator: $Q$ annihilates everything $P$ produces.

\emph{Step 3: range of $P$ is orthogonal to range of $Q$.} Take any $v, w \in V$. We
test the inner product of an arbitrary range vector $Pv$ against an arbitrary leftover
vector $Qw$:
\begin{align*}
  \langle Pv,\; Qw \rangle
    \;&=\; \langle Pv,\; (I-P)w \rangle \\
    \;&=\; \langle (I-P)^\ast Pv,\; w \rangle
        &&\text{(move the operator to the other slot via the adjoint)} \\
    \;&=\; \langle (I-P) Pv,\; w \rangle
        &&\text{(since $(I-P)^\ast = I-P$ by Step 1)} \\
    \;&=\; \langle QP\, v,\; w \rangle \\
    \;&=\; \langle 0,\; w \rangle
        &&\text{(since $QP = 0$ by Step 2)} \\
    \;&=\; 0 .
\end{align*}
Because $v$ and $w$ were arbitrary, \emph{every} vector $Pv$ in the range of $P$ is
orthogonal to \emph{every} leftover $Qw = w - Pw$. In particular, taking $w$ to range
over all of $V$, the leftover $v - Pv$ is orthogonal to the entire range $W$.

\emph{Step 4: closest point (Pythagoras).} Fix $v$ and let $w$ be any point of $W$.
Since $Pv \in W$ and $w \in W$, also $Pv - w \in W$, hence by Step~3 it is orthogonal
to the leftover $v - Pv$. The Pythagorean law for orthogonal vectors
\citep{fourier_rutgers} then gives
\[
  \|v - w\|^2 \;=\; \|\,(v - Pv) + (Pv - w)\,\|^2
              \;=\; \|v - Pv\|^2 + \|Pv - w\|^2
              \;\ge\; \|v - Pv\|^2,
\]
with equality if and only if $\|Pv - w\|^2 = 0$, i.e.\ $w = Pv$. Therefore $Pv$ is the
unique point of $W$ nearest to $v$. This is exactly what ``$P$ projects $v$
orthogonally onto $W$'' means.
\end{proof}

\begin{remark}
The converse is also true and equally short: an orthogonal projection is automatically
self-adjoint \citep{projectors_garrett,wiki_projection}. We will only need the
direction proved above. The single fact you must carry forward is the slogan
\[
  \boxed{\ \text{idempotent} + \text{self-adjoint} \;\Longrightarrow\;
          \text{orthogonal projection (the closest-point map)}.\ }
\]
\end{remark}

\begin{corollary}[Eigenvalues of a projector are $0$ or $1$]
\label{cor:eig01}
If $P^2 = P$ and $Pv = \lambda v$ with $v \neq 0$, then $\lambda v = Pv = P^2 v =
\lambda(Pv) = \lambda^2 v$, so $\lambda = \lambda^2$, forcing $\lambda \in \{0,1\}$
\citep{wiki_projection}. The eigenvectors with $\lambda = 1$ span the range (what $P$
keeps); those with $\lambda = 0$ span the kernel (what $P$ discards).
\end{corollary}

\subsection{Spectral norm: how much an operator can stretch}

We also need one number that measures the largest stretching a linear map can do,
because later it controls how a shift or a perturbation can grow.

\begin{definition}[Spectral / operator norm]
\label{def:specnorm}
For a matrix $W$ the \emph{spectral norm} (also operator norm) is the largest factor by
which $W$ can lengthen a unit vector:
\[
  \|W\|_2 \;=\; \max_{\|z\| = 1} \|Wz\| .
\]
Equivalently $\|W\|_2 = \sigma_{\max}(W)$, the largest \emph{singular value} of $W$,
which equals $\sqrt{\lambda_{\max}(W^\ast W)}$, the square root of the largest
eigenvalue of $W^\ast W$.
\end{definition}

The equality $\|W\|_2 = \sqrt{\lambda_{\max}(W^\ast W)}$ comes from maximizing the
\emph{Rayleigh quotient}: $\|Wz\|^2 = \langle Wz, Wz\rangle = \langle W^\ast W z,
z\rangle$, and over unit vectors this is largest exactly at the top eigenvector of the
self-adjoint matrix $W^\ast W$, where it equals $\lambda_{\max}(W^\ast W)$.

\begin{example}
For $W = \left[\begin{smallmatrix} 3 & 0 \\ 0 & 1 \end{smallmatrix}\right]$ we have
$W^\ast W = \left[\begin{smallmatrix} 9 & 0 \\ 0 & 1 \end{smallmatrix}\right]$, whose
eigenvalues are $9$ and $1$, so $\|W\|_2 = \sqrt 9 = 3$. A useful sanity check on the
whole section: an \emph{orthogonal} projection never stretches, $\|P\|_2 \le 1$,
because $\|Pv\| \le \|v\|$ is just the inequality $\|v-Pv\|^2 \ge 0$ rearranged through
the Pythagoras identity of Step~4.
\end{example}

\begin{problem}
\label{prob:proj}
Build $P = a a^\top / (a^\top a)$ for $a = (2,1)^\top$, so $P = \tfrac15
\left[\begin{smallmatrix} 4 & 2 \\ 2 & 1 \end{smallmatrix}\right]$. Verify $P^2 = P$
and $P = P^\top$, and check that $P$ annihilates $(-1,2)^\top$ (which is orthogonal to
$a$). Then take the oblique $P' = \left[\begin{smallmatrix} 0 & 0 \\ 1 & 1
\end{smallmatrix}\right]$: confirm $P'^2 = P'$ but $P' \neq P'^\top$, find its range
and kernel, and observe they are not perpendicular. The moral: Theorem~\ref{thm:orthoproj}'s
two hypotheses are both load-bearing.
\end{problem}


% =====================================================================
\section{The discrete Fourier transform and the structure of shifts}
\label{sec:dft}

A previous version of this primer asserted, with no explanation, that ``shift operators
are diagonalized by the DFT.'' That sentence is true and it is the keystone of the
whole subject, but it is useless until every word in it is defined and the claim is
proved. This section does that, building up from scratch, following the
``discover the transform, don't postulate it'' approach of Bamieh
\citep{fourier_bamieh}, with the concrete computations of the MIT~18.06 notes
\citep{fourier_mit1806} and the rigorous theorem--proof spine of the Rutgers notes
\citep{fourier_rutgers}.

\paragraph{Notation, fixed once.}
We work with vectors of length $N$, indexed $0,1,\dots,N-1$, and read every index
\emph{modulo $N$} (so index $-1$ means $N-1$, index $N$ means $0$, and so on; we think
of the entries as sitting on a circle). The basic constant is the \emph{primitive
$N$-th root of unity}
\[
  \omega \;=\; e^{-2\pi i / N},
\]
chosen with a minus sign to match the engineering / forward DFT convention used throughout this
primer. It satisfies $\omega^N = 1$ but $\omega^m \neq 1$ for $0 < m < N$; its powers
$1, \omega, \omega^2, \dots, \omega^{N-1}$ are the $N$ solutions of $z^N = 1$, equally
spaced points on the unit circle. The \emph{discrete Fourier transform} (DFT) matrix
$F$ has entries
\[
  F_{jk} \;=\; \omega^{jk}, \qquad 0 \le j,k \le N-1,
  \qquad\text{and we write } \widehat{x} = \mathrm{DFT}(x) = F x .
\]
Two facts about $\omega$ we will use constantly: $\overline{\omega} = e^{+2\pi i/N} =
\omega^{-1}$ (conjugating flips the sign of the exponent), and the \emph{finite
geometric series}
\begin{equation}
\label{eq:geomsum}
  \sum_{m=0}^{N-1} \omega^{rm} \;=\;
  \begin{cases}
    N, & \text{if } r \equiv 0 \pmod N, \\[2pt]
    \dfrac{1 - \omega^{rN}}{1 - \omega^{r}} = \dfrac{1-1}{1-\omega^r} = 0, &
    \text{otherwise,}
  \end{cases}
\end{equation}
where the first case is ``add $N$ copies of $1$'' and the second uses $\omega^{rN} =
(\omega^N)^r = 1$ while $\omega^r \neq 1$. Equation~\eqref{eq:geomsum} is the engine
behind both Parseval and the diagonalization.

\subsection{What ``diagonalize'' means, and the shift operator}

\begin{definition}[Eigenvector, eigenvalue, diagonalize]
\label{def:diag}
A nonzero vector $v$ is an \emph{eigenvector} of a matrix $A$ with \emph{eigenvalue}
$\lambda$ if $A v = \lambda v$: the matrix only stretches $v$, never tilts it. To
\emph{diagonalize} $A$ is to find an invertible matrix $V$ (its columns a full set of
eigenvectors $v_1,\dots,v_N$) so that $V^{-1} A V = D$ is diagonal, with the
eigenvalues $\lambda_1,\dots,\lambda_N$ on the diagonal. In that eigenvector basis the
matrix acts by simple coordinatewise scaling. A family of matrices is
\emph{simultaneously diagonalized} by one $V$ if the same eigenvector basis works for
all of them at once.
\end{definition}

\begin{definition}[Cyclic shift]
\label{def:shift}
The \emph{cyclic (right) shift} $S$ moves every entry one step forward around the
circle:
\[
  (S x)_k \;=\; x_{k-1} \quad (\text{indices mod } N),
  \qquad\text{e.g.}\quad S\,(x_0,x_1,x_2,x_3) = (x_3,x_0,x_1,x_2).
\]
As a matrix $S$ is a \emph{permutation matrix} (a single $1$ in each row and column).
It is the most fundamental object of this whole section; we now find its eigenvectors.
\end{definition}

\subsection{The shift is diagonalized by the DFT}

\begin{theorem}[Eigenvectors of the shift]
\label{thm:shifteig}
Let $h_j$ denote the $j$-th column of the DFT matrix $F$, that is
\[
  h_j \;=\; \bigl(\omega^{0\cdot j},\ \omega^{1\cdot j},\ \omega^{2\cdot j},\ \dots,\
            \omega^{(N-1)j}\bigr)^{\!\top}
        \;=\; \bigl(1,\ \omega^{j},\ \omega^{2j},\ \dots,\ \omega^{(N-1)j}\bigr)^{\!\top}.
\]
Then each $h_j$ is an eigenvector of the cyclic shift $S$:
\[
  S\, h_j \;=\; \omega^{-j}\, h_j \qquad (j = 0,1,\dots,N-1),
\]
so the eigenvalues of $S$ are the $N$ roots of unity $\{\omega^{-j}\} = \{1, \omega^{-1},
\dots, \omega^{-(N-1)}\}$, the same set as $\{1, \omega, \dots, \omega^{N-1}\}$ in a
different order.
\end{theorem}

\begin{proof}
We chase one entry. The $k$-th entry of $h_j$ is $(h_j)_k = \omega^{kj}$. Apply the
definition of the shift, $(S\,h_j)_k = (h_j)_{k-1}$, and compute, for each $k$,
\begin{align*}
  (S\, h_j)_k \;&=\; (h_j)_{k-1} \\
              \;&=\; \omega^{(k-1)j}
                  &&\text{(definition of $h_j$ at index $k-1$)} \\
              \;&=\; \omega^{kj}\,\omega^{-j}
                  &&\text{(law of exponents $\omega^{(k-1)j} = \omega^{kj}\cdot\omega^{-j}$)} \\
              \;&=\; \omega^{-j}\,(h_j)_k .
\end{align*}
This holds for every entry $k$ (the wrap-around at $k=0$ is automatic because the
exponent is read mod $N$ and $\omega^{N} = 1$). Therefore $S\,h_j = \omega^{-j} h_j$ as
vectors: $h_j$ is an eigenvector with eigenvalue $\omega^{-j}$. The $N$ eigenvalues
$\omega^{0}, \omega^{-1}, \dots, \omega^{-(N-1)}$ are distinct, so the $N$ eigenvectors
$h_0,\dots,h_{N-1}$ are linearly independent and form a basis: the columns of $F$
diagonalize $S$.
\end{proof}

\begin{remark}[A note on conventions, so you are never confused later]
With the \emph{forward} convention $\omega = e^{-2\pi i/N}$ adopted here, the shift
eigenvalue attached to column $h_j$ is $\omega^{-j}$. Some textbooks
(e.g.\ \citep{fourier_mit1806,fourier_rutgers}) use the opposite sign $\omega =
e^{+2\pi i/N}$ and obtain $S h_j = \omega^{j} h_j$. The eigenvectors are the same lines,
the eigenvalues are the same set of roots of unity; only the bookkeeping label changes.
We will always state which convention is in force.
\end{remark}

\subsection{Circulant matrices: every shift-invariant linear map}

\begin{definition}[Circulant matrix and circular convolution]
\label{def:circulant}
A matrix $C$ is \emph{circulant} if each column is the cyclic shift of the previous one;
equivalently its entries depend only on $k - l \bmod N$, written $(C_a)_{kl} =
a_{k-l}$, where $a = (a_0,\dots,a_{N-1})$ is its first column. The action of a circulant
is exactly a \emph{circular convolution}: writing $a \conv x$ for the convolution,
\[
  (C_a x)_k \;=\; \sum_{l=0}^{N-1} a_{k-l}\, x_l \;=\; (a \conv x)_k \qquad
  (\text{indices mod } N).
\]
A convolution ``slides one signal across another and sums the overlaps.'' The shift
itself is the circulant whose first column is $(0,1,0,\dots,0)$.
\end{definition}

The reason circulants matter is that they are \emph{precisely} the matrices that
commute with the shift, and ``commutes with the shift'' is the algebraic spelling of
``treats every position the same way'' --- the defining feature of a translation-blind
operation \citep{fourier_bamieh,groups_bronstein}.

\begin{lemma}[Circulants are polynomials in the shift]
\label{lem:polyS}
A matrix $C$ is circulant if and only if it is a polynomial in $S$,
\[
  C \;=\; c_0 I + c_1 S + c_2 S^2 + \dots + c_{N-1} S^{N-1},
\]
where $(c_0,\dots,c_{N-1})$ is the first column of $C$. In particular every circulant
commutes with $S$, and any two circulants commute with each other.
\end{lemma}

\begin{proof}
Note first that $S^p$ is the circulant that shifts by $p$ places, $(S^p x)_k = x_{k-p}$.
Hence a polynomial $\sum_p c_p S^p$ sends $x$ to the vector with $k$-th entry $\sum_p
c_p x_{k-p}$, which is the circular convolution $(c \conv x)_k$ --- a circulant with
first column $c$. Conversely, given a circulant $C$ with first column $c$, the operator
$\sum_p c_p S^p$ has the same first column ($S^p e_0 = e_p$, so $\sum_p c_p S^p$ applied
to $e_0$ returns $c$) and, being shift-invariant, the same every column; so the two
matrices agree. Finally $S$ commutes with any polynomial in $S$, and two polynomials in
$S$ commute with each other because powers of $S$ do
\citep{fourier_rutgers}.
\end{proof}

\begin{theorem}[The DFT simultaneously diagonalizes all circulants;\\
the convolution theorem]
\label{thm:circdiag}
Let $C_a$ be the circulant with first column $a$, equivalently $C_a x = a \conv x$.
Then every column $h_j$ of $F$ is an eigenvector of $C_a$, and the eigenvalue is the
$j$-th DFT coefficient of $a$:
\[
  C_a\, h_j \;=\; \widehat{a}_j\, h_j,
  \qquad\text{where } \widehat{a}_j = (Fa)_j = \sum_{l=0}^{N-1} a_l\, \omega^{jl}.
\]
Consequently the DFT turns convolution into entrywise multiplication --- the
\emph{convolution theorem}:
\[
  \boxed{\ \mathrm{DFT}(a \conv x) \;=\; \mathrm{DFT}(a)\,\odot\,\mathrm{DFT}(x)\ }
  \qquad (\odot = \text{entrywise product}).
\]
\end{theorem}

\begin{proof}
\emph{Eigenvalues.} Write $C_a = \sum_p a_p S^p$ as in Lemma~\ref{lem:polyS}. Apply it
to the shift eigenvector $h_j$ and use $S h_j = \omega^{-j} h_j$
(Theorem~\ref{thm:shifteig}), hence $S^p h_j = \omega^{-pj} h_j$:
\[
  C_a\, h_j \;=\; \sum_{p=0}^{N-1} a_p\, S^p h_j
            \;=\; \sum_{p=0}^{N-1} a_p\, \omega^{-pj}\, h_j
            \;=\; \Bigl(\underbrace{\sum_{p=0}^{N-1} a_p\, \omega^{-pj}}_{=\,(Fa)_{-j}=\widehat a_{-j}}\Bigr)\, h_j .
\]
The bracket is a single number, so $h_j$ is an eigenvector. (With our forward
convention the eigenvalue on column $h_j$ is $\widehat a_{-j}$; the whole \emph{set} of
eigenvalues is $\{\widehat a_0, \dots, \widehat a_{N-1}\}$, just relabelled. The clean
``eigenvalue $=$ DFT of the first column'' statement holds verbatim with the
opposite-sign convention \citep{fourier_bamieh}; we keep ours and note the relabelling.)
Because the $h_j$ form a basis, $C_a$ is diagonal in that basis: this is what
``simultaneously diagonalized'' means, since the same columns of $F$ work for
\emph{every} choice of $a$.

\emph{Convolution theorem (direct index chase, following
\citep{fourier_rutgers,wiki_convolution}).} We compute the $k$-th DFT coefficient of
$a \conv x$ from the definitions:
\begin{align*}
  \mathrm{DFT}(a \conv x)_k
    \;&=\; \sum_{j=0}^{N-1} (a \conv x)_j\, \omega^{kj} \\
    \;&=\; \sum_{j=0}^{N-1} \Bigl(\sum_{l=0}^{N-1} a_{j-l}\, x_l\Bigr)\, \omega^{kj}
        &&\text{(definition of convolution)} \\
    \;&=\; \sum_{l=0}^{N-1} x_l \sum_{j=0}^{N-1} a_{j-l}\, \omega^{kj}
        &&\text{(swap the two finite sums)} \\
    \;&=\; \sum_{l=0}^{N-1} x_l \sum_{m=0}^{N-1} a_{m}\, \omega^{k(m+l)}
        &&\text{(set $m = j - l$; still a full period)} \\
    \;&=\; \sum_{l=0}^{N-1} x_l\, \omega^{kl} \sum_{m=0}^{N-1} a_{m}\, \omega^{km}
        &&\text{(factor $\omega^{k(m+l)} = \omega^{km}\,\omega^{kl}$)} \\
    \;&=\; \Bigl(\sum_{l} x_l \omega^{kl}\Bigr)\Bigl(\sum_{m} a_m \omega^{km}\Bigr)
        \;=\; \widehat{x}_k\,\widehat{a}_k .
\end{align*}
That is exactly the entrywise product $\widehat a \odot \widehat x$ at index $k$. The
substitution $m = j - l$ is legitimate because both $a$ and the exponential are
$N$-periodic in their index, so summing $j$ over one full period is the same as summing
$m$ over one full period.
\end{proof}

\subsection{Parseval: the DFT preserves energy (up to a factor of $N$)}

\begin{theorem}[Parseval / Plancherel]
\label{thm:parseval}
With the forward (unnormalized) DFT, for every vector $x \in \mathbb{C}^N$,
\[
  \|\mathrm{DFT}(x)\|^2 \;=\; N\,\|x\|^2,
  \qquad\text{equivalently}\qquad \|x\|^2 = \tfrac1N\,\|\mathrm{DFT}(x)\|^2 .
\]
(With the \emph{unitary} normalization $\tfrac{1}{\sqrt N}F$ it reads $\|x\| =
\|\mathrm{DFT}(x)\|$ exactly.)
\end{theorem}

\begin{proof}
The proof is the orthogonality of the columns of $F$, which is exactly the geometric
series \eqref{eq:geomsum}. Take two columns $h_j, h_k$ of $F$ and form their inner
product (recall $\langle u,v\rangle = \sum_m u_m \overline{v_m}$, and
$\overline{\omega^{km}} = \omega^{-km}$):
\begin{align*}
  \langle h_j,\, h_k\rangle
    \;&=\; \sum_{m=0}^{N-1} \omega^{jm}\,\overline{\omega^{km}}
        \;=\; \sum_{m=0}^{N-1} \omega^{jm}\,\omega^{-km}
        \;=\; \sum_{m=0}^{N-1} \omega^{(j-k)m}
        \;=\;
        \begin{cases} N, & j = k, \\ 0, & j \neq k, \end{cases}
\end{align*}
where the last step is \eqref{eq:geomsum} with $r = j-k$. In matrix form this says
$F^\ast F = N I$, i.e.\ the columns are orthogonal each of squared length $N$. Now
expand the DFT in this basis. Writing $\widehat x = Fx$,
\[
  \|\widehat x\|^2
    \;=\; \langle Fx,\, Fx\rangle
    \;=\; \langle F^\ast F x,\, x\rangle
    \;=\; \langle N x,\, x\rangle
    \;=\; N\,\|x\|^2 ,
\]
where we moved $F$ to the other slot of the inner product through its adjoint $F^\ast$
and then used $F^\ast F = N I$. Dividing by $N$ gives the second form.
\end{proof}

\subsection{Worked examples at {\boldmath$N=4$} (do these by hand once)}

With $N = 4$ the root of unity is $\omega = e^{-2\pi i/4} = -i$, and the forward DFT
matrix is
\[
  F_4 \;=\;
  \begin{bmatrix}
    1 & 1 & 1 & 1 \\
    1 & -i & -1 & i \\
    1 & -1 & 1 & -1 \\
    1 & i & -1 & -i
  \end{bmatrix}.
\]
(The entry $F_{jk} = (-i)^{jk}$; the MIT~18.06 notes \citep{fourier_mit1806} print the
complex-conjugate version because they use the $+$ sign convention.)

\begin{example}[Convolution two ways]
Let $a = (1,1,0,0)$ and $x = (1,2,3,4)$. Direct circular convolution: $(a\conv x)_k =
\sum_l a_{k-l} x_l$ with $a_0 = a_1 = 1$, $a_2 = a_3 = 0$, so $(a\conv x)_k = x_k +
x_{k-1}$:
\[
  a \conv x = (x_0 + x_3,\ x_1 + x_0,\ x_2 + x_1,\ x_3 + x_2)
            = (1+4,\ 2+1,\ 3+2,\ 4+3) = (5,3,5,7).
\]
Via the transform: $\widehat a = F_4 a = (2,\,1-i,\,0,\,1+i)$ and $\widehat x = F_4 x =
(10,\,-2+2i,\,-2,\,-2-2i)$. Their entrywise product is $\widehat a \odot \widehat x =
(20,\,(1-i)(-2+2i),\,0,\,(1+i)(-2-2i)) = (20,\,4i,\,0,\,-4i)$, hmm let us simplify:
$(1-i)(-2+2i) = -2+2i+2i-2i^2 = -2+4i+2 = 4i$ and similarly the last entry is $-4i$.
Applying the inverse DFT, $\tfrac14 F_4^\ast(20,4i,0,-4i)$, returns $(5,3,5,7)$,
matching the direct computation \citep{fourier_rutgers}. The two routes agree exactly,
which is the convolution theorem at work.
\end{example}

\begin{example}[Parseval check]
For $x = (1,2,-1,0)$ we have $\|x\|^2 = 1 + 4 + 1 + 0 = 6$. Its DFT is $\widehat x = F_4
x = (2,\,2-2i,\,-2,\,2+2i)$, so $\|\widehat x\|^2 = |2|^2 + |2-2i|^2 + |-2|^2 +
|2+2i|^2 = 4 + 8 + 4 + 8 = 24 = 4 \cdot 6 = N\|x\|^2$, exactly as
Theorem~\ref{thm:parseval} predicts.
\end{example}

\begin{problem}
\label{prob:dft}
Show by hand that the $4\times 4$ cyclic shift $S$ has eigenvalues equal to the $4$th
roots of unity $\{1, i, -1, -i\}$, with eigenvectors the columns of $F_4$. Then verify
Parseval for $x = (1,2,-1,0)$ (the example above) and for one vector of your own
choosing. Finally, confirm the convolution theorem on $a = (1,1,0,0)$, $x = (1,2,3,4)$
by reproducing both routes in the previous example.
\end{problem}


% =====================================================================
\section{Symmetry made precise: group actions, orbits, invariance,
and group averaging}
\label{sec:groups}

The DFT story was about \emph{one} symmetry, the cyclic shift. To talk about
shift-invariance and adversarial robustness in general we need the vocabulary that
makes ``a symmetry acting on data'' precise. The right language is that of group
actions, and the right tool is the operation that ``averages a function over a
symmetry'' --- the classical Reynolds operator. We follow the geometric deep learning
treatment of Bronstein, Bruna, Cohen and Veli\v{c}kovi\'c \citep{groups_bronstein} and
the Wikipedia articles on group actions and equivariant maps
\citep{wiki_groupaction,wiki_equivariant}, downgrading their integrals over abstract
groups to finite sums an undergraduate can compute.

\subsection{Groups, actions, orbits}

\begin{definition}[Group]
\label{def:group}
A \emph{group} $(G,\cdot)$ is a set $G$ with a multiplication $\cdot$ that is
associative, has an \emph{identity} element $e$ (with $e g = g e = g$ for all $g$), and
gives every $g$ an \emph{inverse} $g^{-1}$ (with $g g^{-1} = g^{-1} g = e$). The group is
\emph{abelian} (commutative) if $g h = h g$ always. The running example is the
\emph{cyclic group} $\Z_N = \{0,1,\dots,N-1\}$ under addition mod $N$, which is abelian
and is generated by the single element $1$ (every element is a repeated sum of $1$'s).
\end{definition}

\begin{definition}[Group action]
\label{def:action}
A (left) \emph{action} of a group $G$ on a set $X$ is a rule $G \times X \to X$, written
$(g,x) \mapsto g \cdot x$, satisfying two axioms \citep{wiki_groupaction}:
\[
  \text{(identity)}\quad e \cdot x = x, \qquad\qquad
  \text{(compatibility)}\quad g \cdot (h \cdot x) = (g h) \cdot x .
\]
The identity axiom says ``doing nothing does nothing''; the compatibility axiom says
``doing $h$ then $g$ is the same as doing the product $g h$.'' When $X$ is a vector
space and each $g$ acts by a linear (in fact invertible) map, the action is called a
\emph{representation}; we write that map as a matrix $\rho(g)$, and compatibility reads
$\rho(g)\rho(h) = \rho(gh)$.
\end{definition}

\begin{example}[The cyclic shift action, the one we care about]
\label{ex:shiftaction}
Let $G = \Z_N$ act on signals $x \in \mathbb{C}^N$ by shifting: the group element $u$
acts by $\rho(u) = S^u$, so $(\,u \cdot x\,)_k = x_{k - u}$ (shift the signal forward by
$u$). The identity $u = 0$ does nothing; compatibility $S^u S^v = S^{u+v}$ holds because
shifting by $u$ then by $v$ shifts by $u + v$. This is the action that turns
``$\Z_N$-symmetry'' into ``the shift operators of Section~\ref{sec:dft}.''
\end{example}

\begin{definition}[Orbit]
\label{def:orbit}
The \emph{orbit} of a point $x$ is everything reachable from $x$ by the group,
\[
  \Orb(x) \;=\; G \cdot x \;=\; \{\, g \cdot x : g \in G \,\}.
\]
Orbits \emph{partition} $X$: every point is in exactly one orbit, because the relation
``$y = g\cdot x$ for some $g$'' is reflexive ($e$), symmetric ($g^{-1}$) and transitive
(compatibility). An orbit is the set of all configurations that the symmetry regards as
``the same up to the symmetry.''
\end{definition}

\begin{example}[Orbits of the shift on $\R^4$]
\label{ex:orbits4}
Under $\Z_4$ acting by shifts on $\R^4$: the orbit of $(1,0,0,0)$ is the four cyclic
shifts $\{(1,0,0,0),(0,1,0,0),(0,0,1,0),(0,0,0,1)\}$ --- the four standard basis
vectors. The orbit of the constant vector $(1,1,1,1)$ is just itself: shifting it
changes nothing, so it is a \emph{fixed point}. Already you can feel the
shift-invariance theme: a function that should not care about position must give the
same value everywhere along an orbit.
\end{example}

\subsection{Invariant versus equivariant}

This distinction is the single most important conceptual point of the section, and the
old primer never made it.

\begin{definition}[Invariant and equivariant functions]
\label{def:inveq}
Let $G$ act on a domain $X$ and (for equivariance) also on a codomain $Y$. A function
$f : X \to Y$ is
\[
  \textbf{invariant} \quad\text{if}\quad f(g\cdot x) = f(x) \ \ \forall g,
  \qquad\qquad
  \textbf{equivariant} \quad\text{if}\quad f(g\cdot x) = g\cdot f(x) \ \ \forall g .
\]
An invariant function gives the same output no matter how the input is transformed: it
\emph{collapses each orbit to a single value}. An equivariant function lets the
transformation pass through predictably: transform the input and the output transforms
the same way. Invariance is the special case of equivariance where the action on the
codomain is trivial \citep{wiki_equivariant}.
\end{definition}

\begin{example}[The textbook pictures]
\label{ex:areacentroid}
For triangles under the rigid motions of the plane (translations, rotations,
reflections): the \emph{area} is invariant --- moving a triangle does not change its
area. The \emph{centroid} is equivariant --- move the triangle and its centroid moves
to exactly the image of the old centroid \citep{wiki_equivariant}. In statistics, the
sample \emph{mean} is equivariant under rescaling of the data (change units, the mean
changes units the same way), while the \emph{median} is equivariant under \emph{any}
monotone relabelling. In machine learning the canonical pair is: image
\emph{classification} is shift-invariant (a cat is a cat wherever it sits), whereas
\emph{segmentation} is shift-equivariant (shift the image and the predicted mask shifts
identically) \citep{groups_bronstein}.
\end{example}

\begin{remark}[A warning the old primer should have given]
Convolutional layers are shift-\emph{equivariant}, not invariant; this is exactly the
statement $S\,C_a x = C_a\,S x$ for a circulant $C_a$, which you already proved when you
showed circulants commute with $S$ (Lemma~\ref{lem:polyS}). Invariance is only obtained
at the \emph{end}, by a pooling/averaging step that collapses the orbit. Many
signal-processing texts loosely call equivariance ``shift-invariance''; do not be
fooled \citep{groups_bronstein}.
\end{remark}

\subsection{The Reynolds (group-averaging) operator}

How does one \emph{manufacture} an invariant function from an arbitrary one? Average
over the group. This idea is classical (Reynolds; Weyl and Hilbert in invariant theory)
and is the symmetry layer's workhorse \citep{groups_bronstein}.

\begin{definition}[Group average / Reynolds operator]
\label{def:reynolds}
For a finite group $G$ acting on a vector space, define the \emph{group-averaging
operator} $R$ on functions (or, when $G$ acts linearly, on vectors) by
\[
  (R f)(x) \;=\; \frac{1}{|G|} \sum_{g \in G} f(g \cdot x),
  \qquad\text{or for the linear action on vectors:}\quad
  R \;=\; \frac{1}{|G|} \sum_{g \in G} \rho(g).
\]
For a finite group this finite sum is the exact counterpart of Bronstein et al.'s
continuous average $\tfrac{1}{\mu(G)} \int_G f(g\cdot x)\,d\mu(g)$ over the group's
Haar measure \citep{groups_bronstein}: for a finite group the Haar measure just counts
elements.
\end{definition}

\begin{theorem}[The group average is invariant, idempotent, and (for an orthogonal
action) an orthogonal projection onto the invariants]
\label{thm:reynolds}
Let $G$ be finite and act linearly. Then:
\begin{enumerate}[label=\textup{(\roman*)}]
  \item $R$ is \emph{invariant}: $R\,\rho(h) = R$ for every $h \in G$, i.e.\ $(Rf)(h\cdot
        x) = (Rf)(x)$.
  \item $R$ is \emph{idempotent}: $R^2 = R$. Its range is exactly the invariant
        subspace $V_{\mathrm{inv}} = \{v : \rho(g)v = v\ \forall g\}$.
  \item If every $\rho(g)$ is \emph{orthogonal} (or unitary), $\rho(g)^\ast =
        \rho(g)^{-1}$, then $R$ is also self-adjoint, hence by
        Theorem~\ref{thm:orthoproj} the \emph{orthogonal} projection onto
        $V_{\mathrm{inv}}$ --- the closest invariant vector to any given $v$.
\end{enumerate}
\end{theorem}

\begin{proof}
\emph{(i) Invariance, by re-indexing the group.} Fix $h \in G$ and compute, using
compatibility $\rho(g)\rho(h) = \rho(gh)$:
\begin{align*}
  R\,\rho(h)
    \;&=\; \frac{1}{|G|}\sum_{g \in G} \rho(g)\rho(h)
    \;=\; \frac{1}{|G|}\sum_{g \in G} \rho(g h)
    \;=\; \frac{1}{|G|}\sum_{g' \in G} \rho(g')
    \;=\; R .
\end{align*}
The middle step is the only subtlety: as $g$ runs once through all of $G$, the product
$g' = g h$ also runs \emph{once} through all of $G$ (right multiplication by the fixed
$h$ is a bijection of $G$, with inverse right multiplication by $h^{-1}$). So we are
summing the same $|G|$ terms in a different order. Hence $R$ is unchanged by any group
element on its input.

\emph{(ii) Idempotence and range.} First, $R$ fixes every invariant vector: if
$\rho(g)v = v$ for all $g$, then $Rv = \tfrac{1}{|G|}\sum_g \rho(g) v =
\tfrac{1}{|G|}\sum_g v = v$. Second, every output $Rv$ \emph{is} invariant: by part
(i), $\rho(h)\,Rv = R v$ would follow from $\rho(h) R = R$; let us verify that form too,
\[
  \rho(h) R = \frac{1}{|G|}\sum_g \rho(h)\rho(g) = \frac{1}{|G|}\sum_g \rho(hg)
            = \frac{1}{|G|}\sum_{g'} \rho(g') = R ,
\]
the same re-indexing as in (i) but with \emph{left} multiplication. So $Rv \in
V_{\mathrm{inv}}$ for every $v$, and $R$ fixes $V_{\mathrm{inv}}$ pointwise; therefore
$R(Rv) = Rv$, i.e.\ $R^2 = R$, and the range of $R$ is exactly $V_{\mathrm{inv}}$.

\emph{(iii) Self-adjointness for an orthogonal action.} Compute the adjoint of the sum,
using $(\rho(g))^\ast = \rho(g)^{-1} = \rho(g^{-1})$:
\[
  R^\ast = \frac{1}{|G|}\sum_g \rho(g)^\ast
         = \frac{1}{|G|}\sum_g \rho(g^{-1})
         = \frac{1}{|G|}\sum_{g'} \rho(g')
         = R ,
\]
where again $g' = g^{-1}$ ranges over all of $G$ as $g$ does (inversion is a bijection
of $G$). So $R = R^\ast$. Now $R$ is idempotent (part ii) and self-adjoint, so
Theorem~\ref{thm:orthoproj} applies: $R$ is the orthogonal projection onto its range
$V_{\mathrm{inv}}$, the unique nearest invariant vector. This is the rigorous content of
the slogan ``averaging over the symmetry returns the closest symmetric thing.''
\end{proof}

\subsection{The cyclic case, and why linear invariants are too weak}

\begin{example}[Shift-averaging returns the mean]
\label{ex:shiftavg}
Apply Theorem~\ref{thm:reynolds} to $G = \Z_N$ acting on $\mathbb{C}^N$ by shifts,
$\rho(u) = S^u$. The group average is
\[
  (R x)_k = \frac{1}{N}\sum_{u=0}^{N-1} (S^u x)_k
          = \frac{1}{N}\sum_{u=0}^{N-1} x_{k-u}
          = \frac{1}{N}\sum_{j=0}^{N-1} x_j
          = \text{mean}(x),
\]
the constant vector all of whose entries equal the average of $x$. Each $S^u$ is a
permutation matrix, hence orthogonal, so by part~(iii) this $R$ is the orthogonal
projection onto the one-dimensional invariant subspace spanned by $(1,1,\dots,1)$ ---
which is exactly the orbit of fixed points from Example~\ref{ex:orbits4}. Connecting to
Section~\ref{sec:dft}: $R = \tfrac1N\sum_u S^u$ is a circulant, and by
Theorem~\ref{thm:circdiag} its eigenvalues are the DFT of its first column
$\tfrac1N(1,1,\dots,1)$, which is $(1,0,0,\dots,0)$ --- it keeps only the ``zero
frequency'' (the mean) and kills everything else.
\end{example}

The example carries a moral that motivates the entire next section. The \emph{only}
linear shift-invariant feature of a signal is its mean (any other shift-invariant linear
functional must be a multiple of the sum, by the same averaging argument). Bronstein et
al.\ make the same point vividly: for images under translation, a linear invariant can
see nothing but the \emph{average colour} \citep{groups_bronstein}. To get richer
invariants --- ones that can actually tell two signals apart while ignoring their
position --- we must go \emph{nonlinear}. That is precisely what the power spectrum and
bispectrum do.

\begin{problem}
\label{prob:reynolds}
Show that under $\Z_N$ the group average of any vector equals the constant vector whose
entries are the mean; verify directly that $R^2 = R$ on a length-$3$ example; and argue
that the only shift-invariant \emph{linear} functionals $x \mapsto \langle w, x\rangle$
are the multiples of the sum $\sum_k x_k$ (hint: invariance forces $S^\ast w = w$, so
$w$ is itself shift-invariant). Conclude that nontrivial shift-invariants must be
nonlinear.
\end{problem}


% =====================================================================
\section{Shift-invariant descriptors: the power spectrum and the bispectrum}
\label{sec:invariants}

We now build the two nonlinear shift-invariants that the project actually uses. Both
are computed from the DFT, both throw away absolute position, and the difference between
them --- one discards phase, the other keeps relative phase --- is the punchline. The
sources are the MIT~6.011 notes on power spectral density \citep{invariants_mit6011},
de~Buyl's Wiener--Khinchin derivation \citep{invariants_wienerkhinchin}, Kakarala's
introduction to bispectral invariants \citep{invariants_kakarala}, and Chen, Zehni and
Zhao's discrete treatment \citep{invariants_chen}.

\subsection{The DFT shift theorem: shifting only rotates phase}

\begin{lemma}[Shift theorem]
\label{lem:shifttheorem}
Let $R_\tau$ shift a signal by $\tau$, $(R_\tau x)_n = x_{n-\tau}$. Then
\[
  \mathrm{DFT}(R_\tau x)_k \;=\; \widehat{x}_k\, \omega^{k\tau}
                          \;=\; \widehat{x}_k\, e^{-2\pi i k \tau / N}.
\]
A shift leaves the \emph{magnitude} $|\widehat x_k|$ of every Fourier coefficient
untouched and only multiplies it by a \emph{unit-modulus phase} $e^{-2\pi i k\tau/N}$
(a complex number of absolute value $1$).
\end{lemma}

\begin{proof}
Compute from the definition, substituting $m = n - \tau$ and using periodicity:
\begin{align*}
  \mathrm{DFT}(R_\tau x)_k
    \;&=\; \sum_{n=0}^{N-1} (R_\tau x)_n\, \omega^{kn}
    \;=\; \sum_{n=0}^{N-1} x_{n-\tau}\, \omega^{kn}
    \;=\; \sum_{m=0}^{N-1} x_{m}\, \omega^{k(m+\tau)} \\
    \;&=\; \omega^{k\tau}\sum_{m=0}^{N-1} x_m\,\omega^{km}
    \;=\; \omega^{k\tau}\,\widehat x_k. \qedhere
\end{align*}
\end{proof}

Everything below is a consequence of one observation: \emph{position lives entirely in
the phase}. Any quantity built from the $\widehat x_k$ that cannot detect a
phase-rotation $e^{-2\pi i k\tau/N}$ will be shift-invariant.

\subsection{Power spectrum and Wiener--Khinchin}

\begin{definition}[Autocorrelation and power spectrum]
\label{def:powerspec}
The \emph{circular autocorrelation} of $x$ is $r_\tau = \sum_{t=0}^{N-1} x_t\,
x_{t+\tau}$ (for real signals; conjugate the second factor for complex ones), measuring
how much $x$ resembles its own shift by $\tau$. The \emph{power spectrum} is the squared
magnitude of the DFT,
\[
  P_k \;=\; |\widehat x_k|^2 \;=\; \widehat x_k\,\overline{\widehat x_k}.
\]
It is real, nonnegative, and (for real $x$) symmetric \citep{invariants_mit6011}.
\end{definition}

\begin{theorem}[Power spectrum is shift-invariant]
\label{thm:psinv}
For every shift $\tau$, $\ |\mathrm{DFT}(R_\tau x)_k|^2 = |\widehat x_k|^2$.
\end{theorem}

\begin{proof}
Using the shift theorem (Lemma~\ref{lem:shifttheorem}) and $|e^{-2\pi i k\tau/N}| = 1$:
\[
  |\mathrm{DFT}(R_\tau x)_k|^2
    = \bigl|\,\widehat x_k\, e^{-2\pi i k\tau/N}\,\bigr|^2
    = |\widehat x_k|^2 \cdot \bigl|e^{-2\pi i k\tau/N}\bigr|^2
    = |\widehat x_k|^2 \cdot 1
    = |\widehat x_k|^2 .
\]
The phase, which is where the position information sat, is annihilated by the
absolute-value square. Hence the power spectrum is the same for $x$ and every shift of
$x$: it is a shift-invariant descriptor.
\end{proof}

\begin{theorem}[Wiener--Khinchin: power spectrum $=$ transform of autocorrelation]
\label{thm:wk}
The power spectrum is the DFT of the autocorrelation: $\widehat r = P$, equivalently the
autocorrelation is the inverse DFT of the power spectrum.
\end{theorem}

\begin{proof}[Proof skeleton (continuous version, following de~Buyl
\citep{invariants_wienerkhinchin})]
Start from the autocorrelation
\[
  C(\tau) = \lim_{T\to\infty} \tfrac1T \int_0^T x(t)\, x(t+\tau)\,dt
\]
and replace each $x$ by its inverse Fourier transform. After exchanging
the order of integration one is left with a time-average of $e^{i(\omega+\omega')t}$:
that average is $0$ unless $\omega + \omega' = 0$ (the oscillations cancel to order
$1/T$) and equals $1$ when $\omega + \omega' = 0$. Keeping only the surviving
$\omega' = -\omega$ terms and using realness $\widetilde x(-\omega) =
\overline{\widetilde x(\omega)}$ gives
\[
  C(\tau) = \frac{1}{2\pi}\int e^{i\omega\tau}\,\widetilde x(\omega)\,
            \overline{\widetilde x(\omega)}\,d\omega
          = \frac{1}{2\pi}\int e^{i\omega\tau}\,|\widetilde x(\omega)|^2\,d\omega,
\]
i.e.\ the autocorrelation is the inverse transform of $|\widetilde x|^2$, the power
spectrum. The discrete version replaces the integrals by the length-$N$ DFT; the random
(noisy) version is handled by MIT~6.011's periodogram argument
\citep{invariants_mit6011}.
\end{proof}

Why does this matter? Because it shows the power spectrum is a \emph{complete summary of
the second-order structure} of a signal --- but \emph{only} the second order. By
discarding all phase it cannot tell apart two signals whose Fourier magnitudes agree;
Kakarala's Figure~1 shows two genuinely different shapes with identical power spectra,
the power spectrum ``fooled'' \citep{invariants_kakarala}. To recover the lost
discriminative power without giving up shift-invariance, we keep \emph{relative} phase.

\subsection{The bispectrum: a shift-invariant that remembers relative phase}

\begin{definition}[Triple correlation and bispectrum]
\label{def:bispec}
The \emph{triple correlation} is the autocorrelation with one extra shift,
$a_3(s_1,s_2) = \sum_t \overline{x_t}\, x_{t+s_1}\, x_{t+s_2}$. Its DFT is the
\emph{bispectrum}, which in frequency variables is the triple product
\[
  B_{k_1,k_2} \;=\; \widehat x_{k_1}\,\widehat x_{k_2}\,\overline{\widehat x_{k_1+k_2}}
\]
\citep{invariants_kakarala,wiki_bispectrum}. Chen--Zehni--Zhao write the equivalent form
\[
  B_{k_1,k_2} \;=\; \widehat x_{k_1}\,\overline{\widehat x_{k_2}}\,\widehat x_{k_2-k_1}
\]
\citep{invariants_chen}; the two differ only by a relabelling of the frequency axes. Setting
$k_2 = 0$ recovers the power spectrum: $B_{k,0} = \widehat x_k\,\widehat x_0\,
\overline{\widehat x_k} = \widehat x_0\,|\widehat x_k|^2$, so the bispectrum
\emph{contains} the power spectrum \citep{invariants_kakarala}.
\end{definition}

\begin{theorem}[Bispectrum is shift-invariant, by explicit phase cancellation]
\label{thm:bispecinv}
For every shift $\tau$, $\ B^{(R_\tau x)}_{k_1,k_2} = B^{(x)}_{k_1,k_2}$.
\end{theorem}

\begin{proof}
Apply the shift theorem (Lemma~\ref{lem:shifttheorem}) to each of the three Fourier
coefficients in the triple product. Write $\zeta = e^{-2\pi i \tau/N}$, so a shift
multiplies $\widehat x_k$ by $\zeta^{k}$, and \emph{conjugating} $\widehat x_{k_1+k_2}$
multiplies by $\overline{\zeta^{\,k_1+k_2}} = \zeta^{-(k_1+k_2)}$:
\begin{align*}
  B^{(R_\tau x)}_{k_1,k_2}
    \;&=\; \mathrm{DFT}(R_\tau x)_{k_1}\;
           \mathrm{DFT}(R_\tau x)_{k_2}\;
           \overline{\mathrm{DFT}(R_\tau x)_{k_1+k_2}} \\
    \;&=\; \bigl(\zeta^{k_1}\widehat x_{k_1}\bigr)\,
           \bigl(\zeta^{k_2}\widehat x_{k_2}\bigr)\,
           \overline{\bigl(\zeta^{k_1+k_2}\widehat x_{k_1+k_2}\bigr)}
        &&\text{(shift theorem on each factor)} \\
    \;&=\; \zeta^{k_1}\,\zeta^{k_2}\,\zeta^{-(k_1+k_2)}\;
           \widehat x_{k_1}\,\widehat x_{k_2}\,\overline{\widehat x_{k_1+k_2}}
        &&\text{(conjugation flips the phase)} \\
    \;&=\; \zeta^{\,k_1 + k_2 - (k_1 + k_2)}\;
           \widehat x_{k_1}\,\widehat x_{k_2}\,\overline{\widehat x_{k_1+k_2}} \\
    \;&=\; \zeta^{0}\; B^{(x)}_{k_1,k_2}
        \;=\; B^{(x)}_{k_1,k_2} .
\end{align*}
The three shift-phases cancel exactly: $k_1 + k_2 - (k_1+k_2) = 0$, so $\zeta^0 = 1$.
That cancellation is the whole point --- it is why the bispectrum is shift-invariant
\citep{invariants_kakarala,invariants_chen}. Unlike the power spectrum, though, the
bispectrum did \emph{not} have to take an absolute value to achieve invariance: it kept
the \emph{relative} phase among the three coefficients (the part of the phase that a
global shift cannot touch), and that retained phase is exactly the extra information
that lets it separate signals the power spectrum confuses.
\end{proof}

\subsection{Worked examples at {\boldmath$N=4$}}

\begin{example}[A shift leaves $|\mathrm{DFT}|^2$ unchanged]
\label{ex:psn4}
Take $x = (1,0,0,0)$. Then $\widehat x = F_4 x = (1,1,1,1)$, so $P = (1,1,1,1)$. Now
shift to $x' = (0,1,0,0)$: $\widehat{x'} = (1,-i,-1,i)$, and $|\widehat{x'}|^2 =
(1,1,1,1)$ again --- identical power spectrum, even though every phase changed.
\end{example}

\begin{example}[Discrete Wiener--Khinchin by hand]
\label{ex:wkn4}
Take $x = (2,1,0,1)$. Its circular autocorrelation $r_\tau = \sum_t x_t x_{t+\tau}$ is
$r = (6,4,2,4)$ (e.g.\ $r_0 = 4+1+0+1 = 6$). Its DFT is $\widehat r = F_4 r =
(16,4,0,4)$. Independently, $\widehat x = F_4 x = (4,2,0,2)$, so the power spectrum is
$|\widehat x|^2 = (16,4,0,4)$. The two agree: $\widehat r = |\widehat x|^2$, exactly
Theorem~\ref{thm:wk}.
\end{example}

\begin{example}[Same power spectrum, different bispectrum --- the disambiguation]
\label{ex:bispecn4}
Build two real length-$4$ signals
\[
  x_A = (1.25,\ 0.25,\ 0.25,\ 0.25), \qquad
  x_B = (0.75,\ -0.25,\ 0.75,\ 0.75).
\]
Their DFTs are $\widehat{x}_A = (2,1,1,1)$ and $\widehat{x}_B = (2,\,i,\,1,\,-i)$, so
both have the \emph{identical} magnitude spectrum $|\widehat x| = (2,1,1,1)$ and hence
the identical power spectrum $(4,1,1,1)$. The power spectrum cannot tell them apart. The
bispectrum can: with $B_{k_1,k_2} = \widehat x_{k_1}\widehat x_{k_2}
\overline{\widehat x_{k_1+k_2}}$,
\[
  B^{(A)}_{1,1} = \widehat{x}_{A,1}\,\widehat{x}_{A,1}\,\overline{\widehat{x}_{A,2}}
               = 1\cdot 1 \cdot 1 = +1,
  \qquad
  B^{(B)}_{1,1} = \widehat{x}_{B,1}\,\widehat{x}_{B,1}\,\overline{\widehat{x}_{B,2}}
               = i\cdot i \cdot 1 = -1 .
\]
The two differ ($+1$ versus $-1$), entirely because of the relative phase on the first
coefficient. And since both are shift-invariant, this distinction survives every shift:
shifting $x_B$ to $(0.75,0.75,-0.25,0.75)$ leaves $|\widehat x|^2 = (4,1,1,1)$ and
$B_{1,1} = -1$ unchanged. This is the discrete analogue of Kakarala's ``two shapes the
power spectrum confuses but the bispectrum separates'' \citep{invariants_kakarala}.
\end{example}

\begin{problem}
\label{prob:bispec}
For $x = (2,1,0,1)$: (a) verify the discrete Wiener--Khinchin identity $\widehat r =
|\widehat x|^2$ as in Example~\ref{ex:wkn4}; (b) shift $x$ by $2$ and confirm
$|\widehat x|^2$ is unchanged while the individual phases rotate; (c) construct a second
signal with the same $|\widehat x|$ but a flipped phase on one coefficient, and show
that some bispectrum entry $B_{k_1,k_2}$ differs --- so the bispectrum distinguishes
what the power spectrum cannot, even up to shift.
\end{problem}

\paragraph{Where this leaves us.} The four sections fit together as one thread. A
shift-invariant feature is what you get by averaging away the symmetry
(Section~\ref{sec:groups}); doing it \emph{linearly} keeps only the mean, so useful
invariants are nonlinear (Section~\ref{sec:invariants}); the natural nonlinear ones live
in the Fourier domain, where shifting is just a phase rotation (Section~\ref{sec:dft}),
and the whole averaging-is-projection story rests on the one projector theorem of
Section~\ref{sec:projectors}. The power spectrum is the phase-blind invariant; the
bispectrum is the phase-retaining one; both are exactly what the project measures when
it asks whether a network has learned to ignore position.

% =====================================================================
% \citep KEYS USED IN THIS FILE (for the assembler's bibliography):
%   Section 1 (projectors):
%     projectors_garrett   -- Garrett, UMN, "Projectors, projection maps,
%                             orthogonal projections" (central theorem proof)
%     projectors_tamu      -- TAMU MATH423 lecture 36 (numeric examples)
%     wiki_projection      -- Wikipedia, "Projection (linear algebra)"
%                             (oblique counterexample, eigenvalues in {0,1})
%     fourier_rutgers      -- Goodman, Rutgers Math357 (Pythagorean law,
%                             inner-product / Parseval machinery)
%   Section 2 (DFT / circulants):
%     fourier_bamieh       -- Bamieh, "Discovering Transforms" (arXiv:1805.05533)
%     fourier_mit1806      -- Johnson, MIT 18.06 circulant-matrices notes
%     fourier_rutgers      -- Goodman, Rutgers Math357 (diagonalization,
%                             convolution theorem, Parseval proofs)
%     wiki_convolution     -- Wikipedia, "Convolution theorem"
%     groups_bronstein     -- Bronstein-Bruna-Cohen-Velickovic 2021
%   Section 3 (groups / actions / Reynolds):
%     wiki_groupaction     -- Wikipedia, "Group action"
%     wiki_equivariant     -- Wikipedia, "Equivariant map"
%     groups_bronstein     -- Bronstein-Bruna-Cohen-Velickovic 2021 (GDL book)
%       (also attribute classically: Reynolds; Weyl/Hilbert invariant theory;
%        Cohen-Welling 2016; Zaheer et al. 2017 -- prose attributions only)
%   Section 4 (power spectrum / bispectrum):
%     invariants_mit6011        -- MIT 6.011 power spectral density notes
%     invariants_wienerkhinchin -- de Buyl, Wiener-Khinchin derivation
%     invariants_kakarala       -- Kakarala 2012, bispectral Fourier descriptors
%     invariants_chen           -- Chen, Zehni & Zhao 2018 (the file mislabeled
%                                  "bendory"); genuine Bendory et al. is its ref
%     wiki_bispectrum           -- Wikipedia, "Bispectrum"
%       (also attribute classically: Wiener; Khinchin; Einstein 1914;
%        Tukey/Brillinger polyspectra -- prose attributions only)
%
%  FULL \citep KEY LIST:
%    projectors_garrett, projectors_tamu, wiki_projection,
%    fourier_bamieh, fourier_mit1806, fourier_rutgers, wiki_convolution,
%    groups_bronstein, wiki_groupaction, wiki_equivariant,
%    invariants_mit6011, invariants_wienerkhinchin, invariants_kakarala,
%    invariants_chen, wiki_bispectrum
% =====================================================================
